\documentclass{article}
\usepackage{latexml}
\usepackage{alltt}
\usepackage{hyperref}
\usepackage{../sty/latexmldoc}
\input{releases.tex}
\newcommand{\PDFIcon}{\includegraphics{pdf}}

\title{\LaTeXML\ \emph{A \LaTeX\ to XML Converter}}
\begin{document}
\maketitle

%============================================================
The \emph{Current release} is \htmlref{\CurrentVersion}{get}.

In the process of developing the
\href{http://dlmf.nist.gov/}{Digital Library of Mathematical Functions},
we needed a means of transforming
the LaTeX sources of our material into XML which would be used
for further manipulations, rearrangements and construction of the web site.
In particular, a true `Digital Library' should focus on the \emph{semantics}
of the material, and so we should convert the mathematical material into both
content and presentation MathML.
At the time, we found no software suitable to our needs, so we began
development of LaTeXML in-house.  

In brief, \texttt{latexml} is a program, written in Perl, that attempts to
faithfully mimic \TeX's behaviour, but produces XML instead of dvi.
The document model of the target XML makes explicit the model implied
by \LaTeX.
The processing and model are both extensible; you can define
the mapping between \TeX\ constructs and the XML fragments to be created.
A postprocessor, \texttt{latexmlpost} converts this
XML into other formats such as HTML or XHTML, with options
to convert the math into MathML (currently only presentation) or images.

\emph{Caveats}: It isn't finished, there are gaps in the coverage,
particularly in missing implementations of the many useful \LaTeX\ packages.
But is beginning to stabilize and interested parties
are invited to try it out, give feedback and even to help out.


%============================================================
\section{Examples}\label{examples}
At the moment, the best example of \LaTeXML's output is 
the \href{http://dlmf.nist.gov/}{DLMF} itself
--- the preview  will soon contain several chapters.
It is perhaps not a fair example, in that there is much
special-purpose processing to generate cross-referencing and metadata that is not
part of LaTeXML by default. Nevertheless, it gives an indication of the
math conversion and other features possible.

\begin{description}
\item[\href{http://dlmf.nist.gov/}{DLMF}]
   The Digital Library of Mathematical Functions was the
   primary instigator for this project.
\item[\href{manual/}{LaTeXML Manual}]
   The \LaTeXML\ User's manual (\href{manual.pdf}{\PDFIcon}).
\item[\href{examples/tabular/tabular.html}{LaTeX tabular}]
    from the \LaTeX\ manual, p.205.
    (\href{examples/tabular/tabular.tex}{\TeX},
     \href{examples/tabular/tabular.pdf}{\PDFIcon})
\item[These pages] were produced using \LaTeXML, as well.
\end{description}


%============================================================
\section{Get \LaTeXML}\label{get}
\paragraph{Current Release:}\label{download.current}
The current release is \textbf{\CurrentVersion}.
See also the \href{Changes}{Change Log}.


\LaTeXML\ itself is written in `pure' Perl, but requires \TeX,
and several perl modules, some of which depend on C libraries.
It should always be possible to install \LaTeXML\ on any platform
with a sufficiently recent \href{http://www.perl.org/}{Perl},
\ref{install.source}, using CPAN to install the prerequisites.

In most cases, however, it will be preferable to install
a platform-specific, pre-compiled, release, as this often
takes care of installing dependencies and updates,
and allows cleaner uninstall.  See  \ref{install.osnotes} for
platform-specific releases and instructions.

\subsection{Prerequisites}

The following additional Perl modules are required, but often will
be automatically installed by OS-specific procedures:
% %3A%3A for :: ???
\begin{description}
\item[\href{http://search.cpan.org/search?query=Parse::RecDescent&mode=module}{Parse::RecDescent}]
    a useful grammar based parser module.
\item[\href{http://search.cpan.org/search?query=Image::Magick&mode=module}{Image::Magick}]
    provides bindings to the \href{http://www.imagemagick.org/}{ImageMagick} library.
\item[\href{http://search.cpan.org/search?query=XML::LibXML&mode=module}{XML::LibXML}]
    provides bindings to the C library libxml2
    (available from \href{http://www.xmlsoft.org}{xmlsoft}).
    Versions before 1.61 will require the additional \texttt{XML::LibXML::XPathContext} module.
\item[\href{http://search.cpan.org/search?query=XML::LibXSLT&mode=module}{XML::LibXSLT}]
    provides bindings to the C library libxslt
    (from \href{http://www.xmlsoft.org}{xmlsoft},
%  <li><tt><a href="http://search.cpan.org/search?query=XML::LibXML::XPathContext&mode=module">XML::LibXML::XPathContext</a></tt>
%      Depends on the XML::LibXML module.
\item[\href{http://search.cpan.org/search?query=DB_File&mode=module}{DB\_File}]
    usually part of a standard Perl installation, provided
    \href{http://www.sleepycat.com}{BerkeleyDB} is installed.
\item[\href{http://search.cpan.org/search?query=Test::Simple&mode=module}{Test::Simple}]
    usually part of a standard Perl installation.
\end{description}

\subsection{From Source}\label{install.source}
\paragraph{Install Prerequisites}
Assuming that all the required C packages (for libxml2 and libxslt) are already installed on your system,
  use CPAN to install the required perl modules (typically as root):
\begin{alltt}
   perl -MCPAN -e shell
   install Parse::RecDescent, XML::LibXML, XML::LibXSLT, DB_File
   quit
\end{alltt}
The \texttt{Image::Magick} module often installs best by enabling the
perl bindings as part of compiling and installing the main ImageMagick package.

\paragraph{Download} \CurrentTarball
\paragraph{Install} proceeds via the usual Perl module incantations:
\begin{alltt}
   tar zxvf LaTeXML-\CurrentVersion.tar.gz\\
   cd LaTeXML-\CurrentVersion\\
   perl Makefile.PL
   make
   make test
\end{alltt}
and then, as root:
\begin{alltt}
   make install
\end{alltt}
(See \texttt{perl perlmodinstall} for more details, if needed.)

You can specify nonstandard place to install files (and avoid the need to install as root)
by modifying the Makefile creating command, as
\begin{alltt}
   perl Makefile.PL PREFIX=\emph{perldir} TEXMF=\emph{texdir}
\end{alltt}
where \emph{perldir} is where you want the perl related files to go and
\emph{texdir} is where you want the \TeX\ style files to go.

\subsection[OS-Specific Notes]{Operating System Specific Notes}\label{install.osnotes}
With \emph{no} implied endorsement of any of these systems.

\subsubsection[Fedora-based systems]{Fedora-based systems}\label{install.fedora}
For Fedora (and perhaps other similar) distributions, \texttt{yum} will automatically install
the needed prerequisites.
\paragraph{Download} \CurrentFedora
\paragraph{Install} using the command:
\begin{alltt}
   yum --nogpgcheck localinstall LaTeXML-\CurrentVersion-*.ltxml.noarch.rpm
\end{alltt}
If you had previously installed \LaTeXML\ rpm, you may want to explicitly
uninstall it \emph{first}, since the name has changed:
\begin{alltt}
   yum remove perl-LaTeXML
\end{alltt}

\subsubsection[Enterprise systems]{Enterprise-style RPM-based systems (RedHat, Centos)}\label{install.enterprise}
On Red Hat Enterprise Linux and derivatives (Centos, Scientific Linux), \texttt{yum} will also
install prerequisites, but two of them are not available in the current release, 5.
\paragraph{Download Prerequisites} choose from the following according to your architecture:
\begin{description}
\item[32bit]
   \href{releases/perl-XML-LibXML-XPathContext-0.07-1.c5.ltxml.i386.rpm}{perl-XML-LibXML-XPathContext},
   \href{releases/perl-XML-LibXSLT-1.58-1.c5.ltxml.i386.rpm}{perl-XML-LibXSLT}
\item[64bit]
   \href{releases/perl-XML-LibXML-XPathContext-0.07-1.c5.ltxml.x86_64.rpm}{perl-XML-LibXML-XPathContext},
   \href{releases/perl-XML-LibXSLT-1.58-1.c5.ltxml.x86_64.rpm}{perl-XML-LibXSLT}
\item[Source RPM]
    \href{releases/perl-XML-LibXML-XPathContext-0.07-1.c5.ltxml.src.rpm}{perl-XML-LibXML-XPathContext},
    \href{releases/perl-XML-LibXSLT-1.58-1.c5.ltxml.src.rpm}{perl-XML-LibXSLT}
\end{description}
\paragraph{Download} \CurrentCentos
\paragraph{Install} using the command:
\begin{alltt}
   yum --nogpgcheck localinstall LaTeXML-\CurrentVersion-*.ltxml.noarch.rpm \textbackslash\\
       perl-XML-LibXML-XPathContext-0.07-1.*   \textbackslash\\
       perl-XML-LibXSLT-1.58-1.*
\end{alltt}

\subsubsection{Debian-based systems}\label{install.debian}
LaTeXML has been included in the Debian repositories (thanks Atsuhito Kohda);
it should be installable using
\begin{alltt}
  sudo apt-get install latexml
\end{alltt}
This should automatically include any needed dependencies.

Should you wish to install \ref{install.source} or from \ref{download.svn}, the
following command should install the dependencies, after which you may proceed to install LaTeXML.
\begin{alltt}
   sudo apt-get install   \textbackslash\\
      libparse-recdescent-perl \textbackslash\\
      libxml2 libxml-libxml-perl \textbackslash\\
      libxslt1.1 libxml-libxslt-perl  \textbackslash\\
      imagemagick perlmagick
\end{alltt}

%Some \href{http://rhaptos.org/devblog/reedstrm/latexml}{notes} on installation on Debian
%based systems are also available.

\subsubsection{MacOS}\label{install.macos}
LaTeXML has been included in the \href{http://www.macports.org}{MacPorts} repository (thanks Andrew Fernandes);
it should be installable using the command
\begin{alltt}
  sudo port install LaTeXML
\end{alltt}
This should automatically include any needed dependencies.

% \paragraph{For the Adventurous}  As an easy alternative --- if it works ---
% download \CurrentMacOS, save in it's own directory as \texttt{Portfile}
% (without the version number) and, within that directory, run
% \begin{alltt}
%   sudo port install
% \end{alltt}
% This should install \LaTeXML\ and it's all dependencies;
% Otherwise, continue as below.

Should you wish to install \ref{install.source} or from \ref{download.svn}, the
following command should install the dependencies, after which you may proceed to install LaTeXML.
\begin{alltt}
  sudo port install    \textbackslash\\
      p5-xml-libxml p5-xml-libxslt  \textbackslash\\
      p5-parse-recdescent p5-perlmagick
\end{alltt}

\emph{Note} there have been issues reported regarding \verb|DB_File|
not being installed;  Apparently you must install the 
the db `variant' of perl, rather than the gdbm variant;
that is, you must run \verb|sudo port install perl +db|
(possibly after uninstalling perl first?).

\subsubsection{Windows}\label{install.windows}
Many of the prerequisite packages are available in
ppm form (for \href{http://activestate.com/}{ActivePerl} (I think).
I will attempt to generate a ppm for LaTeXML in the near
future, depending on user interest --- and assistance!

There  are reports that LaTeXML runs under
\href{http://strawberryperl.com}{Strawberry Perl}.
The installation runs as follows:

\paragraph{Prerequisites}
\begin{itemize}
\item Install  Strawberry Perl\\
  (This provides \texttt{perl} and \texttt{cpan} commands from the command prompt
   and a handy \texttt{dmake} command to substitute for the Unix \texttt{make}.) 
\item Install \texttt{Parse::RecDescent} by typing the following in the Command Prompt
\begin{alltt}
  cpan -i Parse::RecDescent
\end{alltt}
\item Install an  ImageMagick binary (or from source, if preferred)
\item Install  Ghostscript (otherwise CPAN's Image::Magick won't install right)
\item Try to install Image::Magick via (command prompt)
\begin{alltt}
  cpan -i Image::Magick
\end{alltt}
 If the install fails with a single error on the write.t test, don't worry,
it is a well known error on Win32 systems that is irrelevant.
In that case, proceed to force install within the cpan \emph{shell}
\begin{alltt}
  cpan> force install Image::Magick
\end{alltt}
\end{itemize}
\paragraph{LaTeXML install}
Proceed with the typical install from source, after downloading and unpacking the archive, or checking out the SVN (use TortoiseSVN if in doubt how to access it):
\begin{alltt}
 cd LaTeXML-0.7.0
 perl Makefile.PL
 dmake
 dmake test
 dmake install
\end{alltt}

\subsection{SVN Repository}\label{download.svn}
The most current code can be obtained from the svn repository.
It is browsable at \url{https://svn.mathweb.org/repos/LaTeXML}.
Anonymous checkout is available using the command:
\begin{alltt}
  svn co https://svn.mathweb.org/repos/LaTeXML
\end{alltt}
After checkout, you can continue the installation
as described in \ref{install.source}.

\subsection{License}\label{license}
As this software was developed as part of work done by the
United States Government, it is not subject to copyright,
and is in the public domain.
Note that according to
\href{http://www.gnu.org/licences/license-list.html#PublicDomain}{Gnu.org}
public domain is compatible with GPL.

\subsection{Archived Releases:}\label{archive}
\AllReleases.

%============================================================
\section{Documentation}\label{docs}
If you're lucky, all that should be needed to convert
a \TeX\ file, \textit{mydoc}\texttt{.tex} to XML, and
then to XHTML+MathML would be:
\begin{alltt}
   latexml --dest=\textit{mydoc}.xml \textit{mydoc}
   latexmlpost -dest=\textit{somewhere/mydoc}.xhtml \textit{mydoc}.xml
\end{alltt}
This will carry out a default transformation into XHTML+MathML.  If you
give the destination extension with html, it will generate HTML+images.

If you're not so lucky, or want to get fancy, well \ldots dig deeper:
\begin{description}
\item[\href{manual/index.xhtml}{LaTeXML Manual}]
    Overview of LaTeXML (\href{manual.pdf}{\PDFIcon}).
\item[\href{manual/commands/latexml.xhtml}{\texttt{latexml}}]
    describes the \texttt{latexml} command.
\item[\href{manual/commands/latexmlpost.xhtml}{\texttt{latexmlpost} command}]
   describes the \texttt{latexmlpost} command for postprocessing.
\end{description}

% Possibly, eventually, want to expose:
%   http://www.mathweb.org/wiki/????
% But, it doesn't have anything in it yet.

%============================================================
\section{Contacts \& Support}\label{contact}

\paragraph{Mailing List}
There is a low-volume mailing list for questions, support and comments.
See \href{http://lists.jacobs-university.de/mailman/listinfo/project-latexml}{\texttt{latexml-project}} for subscription information.


\paragraph{Bug-Tracker}
  There is a Trac bug-tracking system for reporting bugs, or checking the
  status of previously reported bugs at
 \href{https://trac.mathweb.org/LaTeXML/}{Bug-Tracker}.


\paragraph{Thanks} to our friends at
the \href{http://kwarc.info}{KWARC Research Group}
for hosting the mailing list, Trac system and svn repository,
as well as general moral support.
Thanks also to \href{http://nist.gov/sima}{Systems Integration for Manufacturing Applications}
for funding portions of the research and development.

\paragraph{Author} \href{mailto:bruce.miller@nist.gov}{Bruce Miller}.
%============================================================
\section{Privacy Notice}\label{privacy}
We adhere to \href{http://www.nist.gov/public_affairs/privacy.htm}{NIST's Privacy, Security and Accessibility Policy}.
%============================================================

\end{document}
